% BOOK-FILE-ID: f-ce46647ce2fe
% 既存の演習を、読む目的に応じて選ぶための案内。
\subsection*{演習の進め方}\label{節：演習の読む経路}
まず例題の解き方を使えるか確かめたい場合は,次の12題に取り組むとよい.その後は他の問題も解き,式の特徴から解法を選ぶ練習をしよう.

\paragraphbf{基本パターン}
\begin{itemize}
  \item \bookproblemref{practice-basic}{10}{標準問10}：等差型.
  \item \bookproblemref{practice-basic}{5}{標準問5}：等比型.
  \item \bookproblemref{practice-basic}{11}{標準問11}：\bookterm{difference}{階差}を書き下して足し合わせる.
  \item \bookproblemref{practice-basic}{3}{標準問3}：階乗で割って等比型にする.
\end{itemize}
\paragraphbf{置き換える式を選ぶ}
\begin{itemize}
  \item \bookproblemref{practice-basic}{9}{標準問9}：\bookterm{fixed}{不動点}からの差.
  \item \bookproblemref{practice-basic}{4}{標準問4}：元の漸化式を満たす一次式との差.
  \item \bookproblemref{practice-basic}{8}{標準問8}：\bookterm{logarithm}{対数}で指数を取り出す.
  \item \bookproblemref{practice-basic}{1}{標準問1}：逆数を使う.
  \item \bookproblemref{practice-basic}{26}{標準問26}：2つの\bookterm{fixed}{不動点}からの差の比.
\end{itemize}
\paragraphbf{対象を組み替える}
\begin{itemize}
  \item \bookproblemref{practice-basic}{6}{標準問6}：隣接3項間を等比型へ分ける.
  \item \bookproblemref{practice-basic}{2}{標準問2}：総和の式をずらして引く.
  \item \bookproblemref{practice-basic}{7}{標準問7}：連立する量の和と差.
\end{itemize}
係数や見た目が変わった場合の確認には,標準の
\bookproblemref{practice-basic}{17}{問17},
\bookproblemref{practice-basic}{20}{問20},
\bookproblemref{practice-basic}{23}{問23},
\bookproblemref{practice-basic}{34}{問34}を使う.
変形を組み合わせる段階では,応用の
\bookproblemref{practice-advanced}{1}{問1},
\bookproblemref{practice-advanced}{7}{問7},
\bookproblemref{practice-advanced}{9}{問9}へ進める.

\paragraphbf{漸化式を立てる}
応用に慣れたら,\,\sectionref*{節：発展入門}で,どの量を数列として置くかを自分で決める練習をしよう.
\bookproblemref{practice-bridge}{1}{発展入門問1}では増える領域の個数,
\bookproblemref{practice-bridge}{2}{問2}では偶奇の移り変わり,
\bookproblemref{practice-bridge}{3}{問3}では短い文字列との対応から式を作る.
\bookproblemref{practice-bridge}{7}{問7}では,連立する量の和と差の意味も言葉で説明してみよう.

\paragraphbf{解法を比較する}
\begin{itemize}
  \item \bookproblemref{practice-basic}{4}{標準問4}と\bookproblemref{practice-basic}{24}{問24}：引く量を探す方法と,\bookterm{difference}{階差}を取る方法や両辺を割る方法は,何を消しているか.
  \item \bookproblemref{practice-basic}{19}{標準問19}と\bookproblemref{practice-basic}{42}{問42}：積・両辺を割る方法・\bookterm{logarithm}{対数}で,条件と計算量はどう違うか.
  \item \bookproblemref{practice-basic}{23}{標準問23}と\bookproblemref{practice-basic}{25}{問25}：特性方程式の重解と,\bookterm{difference}{階差}を重ねる方法はどう結び付くか.
\end{itemize}
短さに加え,初項・係数の変更への対応と例外の見え方も比較する.

\paragraphbf{背景を学ぶ}
\begin{itemize}
  \item \bookproblemref{practice-bridge}{6}{発展入門問6}で箱への分配から和を作った後,\bookproblemref{practice-entrance}{3}{発展問3}へ進める.この総和は,\,\sectionref*{節：演習を別の視点から}で\bookterm{generating}{母関数}の積と比べられる.
  \item \bookproblemref{practice-bridge}{5}{発展入門問5}では位置ベクトルで中点を調べ,\bookproblemref{practice-entrance}{4}{発展問4}では辺上の分点へ広げる.図形の変化は,複素数による回転や\bookterm{matrix}{行列}の積として見直せる.
  \item \bookproblemref{practice-bridge}{8}{発展入門問8}から\bookproblemref{practice-entrance}{6}{発展問6}へは,\bookterm{general}{一般項}に限らず減り方を調べる方法を広げる.\bookproblemref{practice-bridge}{9}{発展入門問9}から\bookproblemref{practice-entrance}{9}{発展問9}へは,整数の関係を保つ漸化式を比較できる.
\end{itemize}

\paragraphbf{操作をつなぐ}
\sectionref*{節：解法のつながり}では,両辺を割る方法と引き算,\bookterm{invariant}{不変量}と\bookterm{linear}{線形}化,
2方向に広がる数表,一周期分の合成,ペル解と近似,二項係数と和を扱う.
本文の方法を確かめるなら問1・2・4・5を選べる.
数表の二方向への広がりは\sectionref*{節：二方向の階差}と問3,
近似分数と和の別の解釈は\chapterref*{章：解法の背景と広がり}と問6・7を往復して読む.
共通する操作を,異なる対象へ使う演習である.
